% NN-G036 further uniform refinement; root whole-proof PASS 2026-11-22T00:06:53Z.
\paragraph{The entire remaining mass-ten gluing branch.}
The reviewed finite projection leaves four mass-ten base rows:
\[
 \begin{aligned}
 &g_1=80,\quad g_2=120,\quad d_1=8,\quad U_1=D,\\
 &(d_2,u_2)\in\{(6,12),(10,10),(14,8),(18,6)\}.
 \end{aligned}
\]
Every allowed auxiliary increment in these rows has
\[
 a_1=8,\qquad v_1=26,\qquad v_2\in\{56,58\}.
\]
Equality with the first auxiliary's absolute trace mass
therefore fixes the entire polynomial, not only its mass:
\[
 A_1=D(D-2)(D-3)
   =2+q+q^{-1}+q^2+q^{-2}+q^3+q^{-3}.                 \tag{G036.T1}
\]
We exclude all four rows, with arbitrary support locations.

The integral $U_1=D$ argument gives
\[
 \begin{aligned}
 G_2&=DWC,\qquad G_1=DH,\qquad H=\delta_1+DC,\\
 C&\equiv D-1\pmod2,\qquad H\equiv0\pmod2,\\
 C(4)&=3,\qquad H(4)=20.
 \end{aligned}                                                   \tag{G036.T2}
\]
The mass-eight cell is necessarily
$\delta_1=\psi+2R$, where $R$ is either two or $T_k$.
In particular $\delta_1(0)$ is four or minus four.

First suppose $W$ is irreducible. If it did not divide
$H$, it would be coprime to $G_1=DH$, because $W(0)\ne0$.
The second auxiliary identity
$V_2G_1=W(DG_2-A_2)$ would force $W\mid V_2$ integrally.
Evaluation at four would then give $10\mid v_2$, contrary
to $v_2=56$ or $58$. Thus $W\mid H$.

In the exceptional case $W=(D-2)B$ of (G036.2),
$B$ is primitive irreducible, $B(4)=5$,
$B\equiv1\pmod2$, and $B(0)=-3$.
If $B\nmid H$, the same auxiliary identity gives
$B\mid V_2$, contradicting $5\nmid v_2$.
Therefore $H=2BK$ for an integral polynomial $K$;
the factor two follows by reducing $H$ modulo two.
But then $H(0)$ is divisible by six, whereas
$H(0)=\delta_1(0)=\pm4$. The exceptional case is impossible.

We are left with the irreducible case, where
\[
 H=2WK,\qquad K\in\mathbb Z[D],\qquad K(4)=1.
\]
Here the integral quotient by $W$ is even because
$\overline W=D$ and $\overline H=0$ in the domain
$\mathbb F_2[D]$. If $W(0)=6$, the identity
$H(0)=2W(0)K(0)=\pm4$ is impossible. Hence
\[
 W=D+2+2T_j,\quad j\ {\rm odd},\qquad W(0)=-2,
 \qquad K(0)\in\{1,-1\}.
\]
Every integer polynomial satisfies $K(4)\equiv K(0)\pmod4$.
Since $K(4)=1$, this forces $K(0)=1$.
Consequently $\delta_1(0)=H(0)=-4$, so
$R=T_k$ with $k$ odd. At $D=2$ the Chebyshev
evaluations are therefore
\[
 W(2)=4,\qquad \delta_1(2)=\psi(2)+2T_k(2)=-2.
                                                                  \tag{G036.T3}
\]

The first auxiliary identity, after cancelling $W$
and then $D$, is
\[
 C V_1=2DWK-A_1/D.
\]
By (G036.T1), evaluation at two gives
\[
                  C(2)V_1(2)=16K(2).               \tag{G036.T4}
\]
The integer $C(2)$ is odd by (G036.T2).
Therefore $16$ divides $V_1(2)$.
The coefficientwise lower bound for $V_1$ is the
effective character $V_{\min}=6+5(q+q^{-1})$,
of mass sixteen. Its parity and mass twenty-six imply
$V_1=V_{\min}+2P$ for an effective symmetric integral
character $P$ of mass five. At $D=2$, equivalently
$q=i$, the triangle inequality gives
$|P(i)|\le5$ and hence
\[
                         -4\le V_1(2)\le16.
\]
Its only multiples of sixteen in this interval are
zero and sixteen.
If $V_1(2)=0$, (G036.T4) gives $K(2)=0$.
The monic polynomial $D-2$ would divide $K$ integrally,
so its value two at four would divide $K(4)=1$,
which is impossible.
Thus $V_1(2)=16$, and (G036.T4) gives $K(2)=C(2)$.
But (G036.T2)--(G036.T3) also give
\[
 C(2)=\frac{2W(2)K(2)-\delta_1(2)}2
      =4K(2)+1.
\]
This would require $-3K(2)=1$, contradicting integrality.
Every mass-ten competitor is excluded.
\hfill$\square$

\paragraph{Meaning of this refinement.}
The closed-table optimum is consequently attained at
gluing dimension eight unless a strictly smaller object
occurs in the remaining dimension-twelve stratum.
The proof uses restriction of integral torus characters
at $q=-1$ and $q=i$, together with ordinary dimensions
at $q=1$; it assumes no Weil-weight estimate and no
effectivity of $C,H/W$, or $K$.
Its geometric objects and source dictionary remain the
actual composing graded sign representations of NN-G010,
SGA~3, I, 4.7.2--4.7.3, and the supplied kernel-product
transcription at archive cote~2, scan~17.
